\section{Introduction}

Severe household food insecurity is commonly summarized for Nigeria at state level, but a state average cannot show where risk is concentrated within a state. The weighted national estimate from the 2024 Nigeria Demographic and Health Survey is 29.10\%. Decisions directed to local government areas (LGAs) require estimates below state level, together with a clear account of how much those estimates and any resulting area count remain uncertain. The relevant target is not only the prevalence in each LGA, but also the number of LGAs whose prevalence crosses a stated threshold.

Household severity and area status are separate decisions. The household outcome defines whether a respondent meets the survey's severe-food-insecurity definition; the LGA threshold classifies an area according to its prevalence. A household outcome cannot by itself identify whether an LGA crosses that line, and an area estimate is not a direct count of affected households. Keeping these levels distinct prevents a household measure, an area prevalence, and a count of places from being treated as the same quantity.

Spatial models of household food-insecurity items provide one relevant precedent. Chacón-Montalván et al. (2025; \href{https://doi.org/10.1214/25-aoas2072}{DOI}) fit a spatial item-factor model to data from Ipixuna in the Brazilian Amazon and map three latent factor scores derived from 18 modified food-insecurity items as a continuous surface. They also report pointwise exceedance probabilities, \(\Pr\{\theta_k(s)>0\}\), use a 0.7 probability rule to label hotspots, and count the resulting hotspot regions: four in the full data, one in the dry season, and four in the rainy season. These results make clear that exceedance probabilities and hotspot counts are already used in spatial food-insecurity work. The mapped quantity, however, is a latent factor score at a point, not an LGA prevalence. Ipixuna is one municipality, and the finest administrative unit discussed is eight census sectors; the paper does not report estimates for administrative areas.

A related article by Moura, Pereira, and Silva (2025), \emph{A Beta--Beta Prime Model for Rates and their Precision for Small Area Estimation: An Application to Brazilian Food Insecurity Index}, is closed access and could not be read in full; we make no claim about its contents (\href{https://doi.org/10.1093/jssam/smaf002}{DOI}).

Other spatial methods address different outcomes and prediction tasks. Michal et al. (2024; \href{https://doi.org/10.1177/09622802241293776}{DOI}) extend the Besag--York--Mollié (BYM2) model with a scale mixture that allows latent variance to vary across areas, propose two priors for the mixture scale, and report 23 of 160 Rio de Janeiro districts as potential outliers in an analysis of Zika cases. MacBride et al. (2025; \href{https://doi.org/10.1214/24-aoas1969}{DOI}) describe SPAR-Forest, which iteratively refits a random forest with a Bayesian conditional autoregressive (CAR) or spatial autoregressive (SAR) model. Their abstract reports better performance than CAR/SAR models, random forests, and geographical random forests on Scottish property prices and in simulations. These studies provide precedents for heavy-tailed area effects and spatial prediction, but they do not change the estimand considered here: the prevalence of severe food insecurity among households in each Nigerian LGA.

Area-level estimation for food insecurity and nutrition is an active practice in settings where surveys are designed for larger units. Chandra and Verma (2022) estimate food-insecurity prevalence for districts of Uttar Pradesh from a survey intended to be reliable at state level; Yilema et al. (2022) combine survey and census data to estimate undernutrition among children under five in Ethiopian districts; Wulandari, Darmawan and Zamzanah (2025) estimate stunting prevalence at subdistrict level in Indonesia with a hierarchical Bayesian model and conditional autoregressive spatial priors; and Santillana et al. (2023) estimate a survey proportion for every municipality in Mexico with the Fay--Herriot area-level model. These papers establish that borrowing strength across areas for nutrition and food insecurity is established practice. None of them reports an interval for the number of areas above a threshold.

The distinction between a spatial surface and an area prevalence matters for both comparison and decisions. A probability that a latent factor is positive is not the probability that more than a chosen share of households in a local government area is severely food insecure. Likewise, a list of areas with high pointwise probabilities does not give the uncertainty in the number of areas above a prevalence threshold. The available evidence must therefore be assessed against the actual decision quantity rather than against a superficially similar map or hotspot label.

Within Nigeria, published estimates exist at two levels. Cross-sectional surveys report household food insecurity for single local government areas, for example Egor in Edo State (Omuemu, Otasowie and Onyiriuka, 2012) and Akinyele in Oyo State (Kayode and Alabi, 2025). At a wider scale, Jerumeh (2024) maps household food insecurity across the 109 senatorial districts from the 2018/19 Nigeria Living Standard Survey, classifies 28 districts as hotspots and 19 as coldspots, and reports spatial dependence between neighbouring districts in a spatial regression; only its published abstract was available to us, its full text not being open access. The senatorial district is an administrative unit above the local government area. In this work we found no interval for an area estimate and no uncertainty statement for the number of areas above a line. Experience-based food insecurity in Nigeria has also been analysed against agricultural productivity (Villacis, Mayorga and Mishra, 2022), and a reproducible machine-learning pipeline predicts household-level food insecurity from the 2024 survey used here, with validation on cluster-grouped splits (Osung, 2026); we found no area-level prevalence estimates for the local government areas in either. Our claim is therefore narrower than the absence of Nigerian estimates: we found no estimates for the national set of 774 local government areas, and no published interval for the count of areas above a threshold.

This scope does not imply that exceedance probabilities are absent from prior work. Chacón-Montalván et al. report them and use a probability rule to count hotspots. Their threshold is zero on a latent factor scale, not a threshold on prevalence, and their hotspot counts are descriptive, without an interval for the count. Their paper reports model-fit comparisons and Markov chain Monte Carlo (MCMC) diagnostics, but no out-of-sample validation. The gaps we can verify in the closest paper concern the geographic unit, a prevalence threshold, uncertainty for the count, and out-of-sample performance.

The uncertainty in a count is not a technical detail. Let \(\theta_s\) denote the severe-food-insecurity prevalence in LGA \(s\), let \(c\) be a chosen threshold, and define \(I_s(c)=\mathbf{1}\{\theta_s>c\}\). The count of areas above the threshold is \(C(c)=\sum_{s=1}^{S}I_s(c)\), where \(S\) is the number of LGAs. If \(p_s=\Pr\{I_s(c)=1\}\) and \(P_{st}=\Pr\{I_s(c)=1,I_t(c)=1\}\), then
\[
  \operatorname{E}\{C(c)\}=\sum_{s=1}^{S}p_s,\qquad
  \operatorname{Var}\{C(c)\}
  =\sum_{s=1}^{S}p_s(1-p_s)
    +2\sum_{s<t}\left(P_{st}-p_sp_t\right).
\]
The mean is fixed by the area-specific probabilities because expectation is additive. The variance is not: it depends on whether exceedances occur together or in different places, information that the marginal probabilities do not contain. In this application, the standard deviation of the count can range from 0.16 to 74.11 under joint distributions that all have the same marginals. This range is not a confidence interval around one fitted estimate. It describes how much the uncertainty in the count can vary when only the marginal probabilities are known. An interval based on one fitted spatial dependence structure is a different, model-dependent quantity, and area-level prediction checks do not by themselves validate that joint count interval. A number with no interval, or with an interval that has never been checked, is not usable for targeting.

This distinction also affects how spatial validation should be interpreted. Prediction for one unsampled LGA while nearby LGAs remain in the fit is not the same task as prediction for an entire region whose local observations and neighbours are withheld. A validation result should match the intended use, and a pass for individual area estimates should not be presented as evidence that the interval for the joint count is calibrated. The paper therefore treats the marginal bounds, the model-based count distribution, and the spatial validation results as separate pieces of evidence.

Cross-validation is a design question, not a single score. Withheld LGAs may be isolated or grouped, and neighbouring observations may remain available or disappear with them. Those designs represent different prediction settings and answer different questions about how far the estimates can be carried. We therefore define the fold structure and pass rules before comparing predictions, rather than treating one split as a universal test of spatial performance. The aim is to identify which prediction tasks the fitted estimates support and which they do not.

We use established probability results and spatial models to examine these questions in the 2024 Nigerian application. The paper makes four contributions:
\begin{enumerate}
  \item We derive sharp bounds on the variance of the number of LGAs above a prevalence threshold when only the marginal exceedance probabilities are specified.
  \item We set out a spatial cross-validation protocol with pass rules declared in advance and apply it in four designs, distinguishing prediction for one area among sampled neighbours from prediction for a whole region without local observations.
  \item We give a partial-identification statement for LGAs with no sampled cluster, separating what the survey identifies from what must remain unknown.
  \item We apply the analysis to severe household food insecurity in Nigeria and find that the environmental covariates tested carry no measurable signal at LGA level.
\end{enumerate}
The contribution is an application and evaluation of existing ideas, not a new estimator. The count bounds clarify what marginal probabilities alone can establish; the holdout designs assess area prediction under distinct spatial information settings; and the partial-identification statement prevents model predictions for unclustered areas from being mistaken for direct observations. Together, these components keep a fitted map, a count interval, and the survey's limits visible as different claims.

The remainder of the paper is organized as follows. Methods defines the outcome, household and area model, count functional, partial-identification bounds, and spatial validation protocol, and presents the supporting propositions. Results reports the descriptive estimates, area summaries, count uncertainty, bounds for areas without clusters, validation outcomes, and covariate analysis. Discussion considers the interpretation and limits of these results. Conclusions summarize the findings, followed by the Declaration and References.
